\section{Methods}
\label{sec:methods}

\subsection{Data Sources and Processing}
The primary dataset utilized in this study is the lunar crustal magnetic field map generated by \citet{Hood2020} (DOI: 10.17189/1520494), derived from vector magnetometer measurements from the Lunar Prospector and SELENE (Kaguya) missions. This dataset provides global coverage between \SI{65}{\degree}S and \SI{65}{\degree}N at an approximately \SI{0.5}{\degree} grid resolution, normalized to an altitude of \SI{30}{\kilo\meter}. To cover the polar regions critical to the Artemis program, we incorporated supplementary polar magnetic maps (\SI{60}{\degree} to \SI{90}{\degree} latitude) compiled by \citet{Hood2022} (DOI: 10.17189/rk57-g992). 

Data processing was performed utilizing custom scripts in Python 3.11. The PDS4 ASCII tables were parsed using the \texttt{pandas} library, and the total magnetic field magnitude ($|\mathbf{B}|$) was computed from the constituent radial, east, and north vector components.

\subsection{Landing Site Characterization}
A comprehensive database of lunar landing site coordinates was constructed, encompassing the 13 candidate sites for the Artemis III mission, historical Apollo sites (Apollo 11, 12, 14, 15, 16, and 17), and recent uncrewed missions including Chang'e 3, 4, 5, and Chandrayaan-3. For each coordinate pair, the local crustal magnetic field magnitude at \SI{30}{\kilo\meter} altitude was extracted using bilinear interpolation from the gridded PDS datasets. 

To systematically evaluate the biological implications of these varied magnetic environments, we developed a classification scheme for the landing sites based on the local field magnitude:
\begin{itemize}
    \item \textbf{High field:} $|\mathbf{B}| > \SI{5}{\nano\tesla}$
    \item \textbf{Moderate field:} $\SI{1}{\nano\tesla} \le |\mathbf{B}| \le \SI{5}{\nano\tesla}$
    \item \textbf{Low/Null field:} $|\mathbf{B}| < \SI{1}{\nano\tesla}$
\end{itemize}

\subsection{Literature Review Methodology}
To synthesize current knowledge on the biological effects of hypomagnetic fields, a systematic literature review was conducted. The PubMed and Web of Science databases were queried for the period between 2000 and 2026. The search strategy employed combinations of the keyword ``hypomagnetic field'' with biological terms, including ``biology'', ``plant'', ``bacteria'', and ``organism''. Studies were filtered for relevance, prioritizing empirical research conducted in controlled HMF environments ($< \SI{5}{\micro\tesla}$) and major systematic reviews.

\subsection{Software and Availability}
All spatial analyses and interpolations were performed using the \texttt{numpy} and \texttt{scipy} libraries. Geospatial visualizations and global maps were generated using \texttt{matplotlib} and \texttt{cartopy}. The complete analysis pipeline and custom scripts developed for this study are archived and publicly accessible on Zenodo (DOI: 10.5281/zenodo.XXXXXXX).
